\documentclass{article}

\usepackage{amsmath,amssymb}
\usepackage{xurl}
\usepackage[hidelinks]{hyperref}

% Shared commands for notes in docs/paper-gaps/.

\DeclareUnicodeCharacter{2080}{\ifmmode{}_{0}\else\textsubscript{0}\fi}
\DeclareUnicodeCharacter{2081}{\ifmmode{}_{1}\else\textsubscript{1}\fi}
\DeclareUnicodeCharacter{2082}{\ifmmode{}_{2}\else\textsubscript{2}\fi}
\DeclareUnicodeCharacter{2099}{\ifmmode{}_{n}\else\textsubscript{n}\fi}

\newcommand{\Lean}{\textsc{Lean}}
\ExplSyntaxOn
\NewDocumentCommand{\leanid}{m}{
  \ifmmode
    \textsf{\detokenize{#1}}
  \else
    \group_begin:
    \urlstyle{sf}
    \tl_set:Nn \l_tmpa_tl {#1}
    \tl_replace_all:Nnn \l_tmpa_tl {\_} {_}
    \exp_args:NV \path \l_tmpa_tl
    \group_end:
  \fi
}
\ExplSyntaxOff
\providecommand{\tr}{\operatorname{tr}}
\newcommand{\tnleanIssueBase}{https://github.com/LionSR/TNLean/issues/}
\newcommand{\tnleanPrBase}{https://github.com/LionSR/TNLean/pull/}
\newcommand{\tnleanDocBase}{https://sirui-lu.com/TNLean/docs/}
\newcommand{\ghissue}[1]{\href{\tnleanIssueBase#1}{GitHub issue \##1}}
\newcommand{\ghpr}[1]{\href{\tnleanPrBase#1}{PR \##1}}
\newcommand{\doclink}[1]{\href{\tnleanDocBase#1.html}{\path{#1}}}

% Dirac notation and the identity operator, as in blueprint/src/macros/common.tex.
% \providecommand keeps note-local definitions authoritative.
\usepackage{dsfont}
\providecommand{\ket}[1]{|#1\rangle}
\providecommand{\bra}[1]{\langle #1|}
\providecommand{\braket}[2]{\langle #1 | #2 \rangle}
\providecommand{\ketbra}[2]{|#1\rangle\!\langle #2|}
\providecommand{\Id}{\mathds{1}}
\providecommand{\one}{\mathds{1}}
\providecommand{\C}{\mathbb{C}}
\providecommand{\MN}[1]{M_{#1}(\C)}
\providecommand{\MD}{M_D(\C)}

% Verdict marker: \gapnote{<kind>}{<status>}. Kind is the policy
% classification (clarification, local-correction, scope-restriction,
% unfaithful, false-source, open-gap); status is open, wip (elimination
% underway), resolved, or historical. Machine-read by the site index;
% prints a verdict line.
\newcommand{\gapnote}[2]{\par\noindent\textbf{Verdict:} #1 (#2).\par}

\importpaperaux{fbc25-}{2502.20257/main-tnlean}

\title{The Length of the Complementary Truncated Symmetry}
\author{}
\date{2026-10-30}

\begin{document}
\maketitle
\gapnote{local-correction}{resolved}

\section{Printed relation}

The truncated symmetry \(U^L_g\) of \papereqref{fbc25}{eq:truncsym} closes a
string of matrix product unitary tensors with the source triangles \(Y_2^g\)
and \(Y_1^g\). The superscript is its total length, including the two endpoint
sites: the composition rule
\(\bigl(\lambda^L_{h^{-1},g^{-1}}\otimes\lambda^R_{g,h}\bigr)U^{L-2}_gU^L_h
=U^L_{gh}\) and the on-site example
\(U^L_g=u_g^{\otimes(L-1)}\otimes\mathbf 1\) both use this reading
\cite[lines~3672 and 5622]{FrancoRubio2025MPUGauging}. The source then states
\begin{align}
    U^L_gU^\infty_{g^{-1}} & =U^{\infty-L}_{g^{-1}},
    \label{eq:fbc25_complementary_printed}
\end{align}
with \(U^{\infty-L}_{g^{-1}}\) the truncated symmetry supported in the
complement of \(U^L_g\) \cite[lines~2175--2181]{FrancoRubio2025MPUGauging}, and
its twisted form \(U^LU^\infty=\sigma U^{\infty-L}\) at an involutive element
\cite[lines~5344--5348]{FrancoRubio2025MPUGauging}, whose proof is left to the
reader.

\section{Shared endpoint sites}

Place a ring of \(H=n+m+2\) sites in the order \((a,p,b,q)\), with \(p\) of
length \(n\) and \(q\) of length \(m\), and let \(U^L\) have total length
\(L=n+2\) and endpoint sites \(a,b\). On \(p\) the product
\(U^LU^\infty_{g^{-1}}\) is the identity. At \(a\) and \(b\) its outputs are
the endpoint legs of \(U^L\), and in general it does not act trivially there.
In the on-site example with \(U^L_g=u_g^{\otimes(L-1)}\otimes\mathbf 1\), the
product acts by \(u_g^{-1}\) on \(b\), which lies in the support of \(U^L_g\),
and by \(u_g^{-1}\) on every site of \(q\). No operator supported on the
\(H-L=m\) sites of \(q\) equals it unless \(u_g\) is a scalar.

The complementary truncated symmetry therefore has bulk \(q\) and the same two
endpoint sites, traversed from \(b\) to \(a\). Its total length is
\begin{align}
    m+2 & =(H-L)+2,
    \label{eq:fbc25_complementary_length}
\end{align}
two more than the printed \(\infty-L\). At each shared endpoint site the leg of
\(U^L\) and the leg of the complement are matched along the identification of
the two source dimensions \(d_l=d_r\); in the source's construction
\(X_1=(\mathbf 1\otimes T^*)Y_2^\dagger\) of \papereqref{fbc25}{eq:modif_XY} the
curly leg of the first-cut factors is the thick leg of \(Y_2\).

\section{Formal treatment}

The formal relation is
\begin{align}
    \bigl(U^{n+2}\otimes I_q\bigr)\,\mathcal U^{(n+m+2)}
     & =\sigma\,\bigl(I_p\otimes U^{m+2}_{(b,q,a)}\bigr)
    \label{eq:fbc25_complementary_formal}
\end{align}
for the inverse-compatible source factors, at every pair of bulk lengths, with
the endpoint legs matched as above.\footnote{Formal declarations:
    \leanid{MPOTensor.SourceFactors.complementaryTruncatedSymmetry} and
    \leanid{MPOTensor.inverseCompatibleSourceFactors_truncatedSymmetry_mul_groupedMPO}
    in \doclink{TNLean/MPS/MPU/InverseCompatiblePleasantProperties}.}
No hypothesis is added; only the length of the complementary truncated
symmetry is corrected.

\textbf{Verdict: resolved local fix.}
The printed complementary length \(\infty-L\) is two short. The complementary
truncated symmetry shares both endpoint sites with \(U^L\), and the corrected
relation (\ref{eq:fbc25_complementary_formal}) is proved.

\bibliographystyle{alpha}
\bibliography{references}

\end{document}
